\chapter{Complexity — Examples}
%%% generated by the blueprint pipeline from TCSlib.Complexity.ClassP.Examples
%%%%%%%%%%%%%%%%%%%%%%%%%%%%%%%%%%%%%%%%%%%%%%%%%%%%%%%%%%%%%%%%%%%%%%%%%%%%%%%
\section{Overview}

This module is the sanity check that the Turing-machine model and the class
definitions are usable: it defines the language of binary palindromes, constructs
an explicit three-state one-work-tape machine that decides it, and proves that the
language lies in $\mathrm{DTIME}(n+1)$ and hence in $\mathsf{P}$
(Arora--Barak, Examples 1.1 and 1.4).

%%%%%%%%%%%%%%%%%%%%%%%%%%%%%%%%%%%%%%%%%%%%%%%%%%%%%%%%%%%%%%%%%%%%%%%%%%%%%%%
\section{Declarations}

\begin{definition}[The language of binary palindromes]
\lean{Complexity.PAL}
\label{Complexity.PAL}
\leanok
The language of binary palindromes: the set of finite words
$x \in \{0,1\}^*$ that equal their own reversal, $x^{R} = x$.
\end{definition}

\begin{definition}[The three-state palindrome machine]
\lean{Complexity.palTM}
\label{Complexity.palTM}
\leanok
\uses{Turing.FinTM}
A concrete Turing machine over the binary alphabet with one work tape and three
control states $\{0,1,2\}$, given by an explicit transition table. State $0$
(\emph{copy}) copies each input bit to the work tape while moving both heads
right; upon reading the blank past the input it switches to state $1$
(\emph{rewind}), which moves the input head back left to the left boundary; state
$2$ (\emph{test}) then moves the input head right and the work head left in
unison, comparing the two scanned symbols, emitting false and halting on any
mismatch and emitting true when the input is exhausted with all pairs matched.
\end{definition}

\begin{definition}[Partially copied work tape]
\lean{Complexity.palTape}
\label{Complexity.palTape}
\leanok
Given a binary word $x$, a time bound $t \in \mathbb{N}$, and an integer cell
index $z$, this is the partial tape content whose cell $z$ holds the bit $x_z$
when $0 \le z < t$ (and $z$ is within $x$), and is blank otherwise. It describes
a work tape onto which exactly the first $t$ bits of $x$ have been copied.
\end{definition}

\begin{definition}[Canonical machine configurations]
\lean{Complexity.palCfg}
\label{Complexity.palCfg}
\leanok
\uses{Complexity.palTape, Turing.Cfg}
For a binary word $x$ of length $n$, a control state $q \in \{0,1,2\}$, an input
head position $p \in \{0,\dots,n+1\}$ (cells $1$ through $n$ hold the input bits,
cells $0$ and $n+1$ are blank boundary cells), an integer work head position $w$,
and a time bound $t$, this is the machine configuration in live state $q$ with
input head at $p$, work tape holding exactly the first $t$ bits of $x$ in cells
$0,\dots,t-1$ (all other cells blank), work head at $w$, and empty output. These
are the configurations appearing in the three phase invariants.
\end{definition}

\begin{lemma}[Writing the next bit extends the copied prefix]
\lean{Complexity.palTape_write}
\label{Complexity.palTape_write}
\leanok
\uses{Complexity.palTape}
\difficulty{4}
Let $x$ be a binary word and $t < |x|$. Updating the partial tape that holds the
first $t$ bits of $x$ (in cells $0,\dots,t-1$, blank elsewhere) by writing the
bit $x_t$ into cell $t$ yields exactly the partial tape holding the first $t+1$
bits of $x$.
\end{lemma}

\begin{lemma}[Single copy-phase transition]
\lean{Complexity.pal_copy_step}
\label{Complexity.pal_copy_step}
\leanok
\uses{Complexity.palCfg, Complexity.palTM, Complexity.palTape_write, Turing.Action.apply, Turing.Cfg.inputSymbol, Turing.MultiTapeTM.step, Turing.inputSymbolInner, Turing.moveInputPos_pos_of_ne_right}
\difficulty{4}
Let $x$ be a binary word and $t < |x|$. From the configuration of the three-state
palindrome machine on input $x$ in copy state with input head at cell $t+1$
(scanning $x_t$), the first $t$ bits copied to the work tape, and work head at
cell $t$, one transition writes $x_t$ at the work head and advances both heads:
it yields the copy-state configuration with input head at $t+2$, the first $t+1$
bits copied, and work head at $t+1$.
\end{lemma}

\begin{lemma}[Copy-phase invariant]
\lean{Complexity.pal_copy}
\label{Complexity.pal_copy}
\leanok
\uses{Complexity.palCfg, Complexity.palTM, Complexity.palTape, Complexity.pal_copy_step, Turing.Cfg.init, Turing.MultiTapeTM.initCfg, Turing.MultiTapeTM.runFrom, Turing.MultiTapeTM.runFrom_succ_eq_step', Turing.MultiTapeTM.runFrom_zero}
\difficulty{4}
Let $x$ be a binary word. For every $t \le |x|$, running the three-state
palindrome machine for $t$ steps from its initial configuration on input $x$
reaches the copy-state configuration with input head at cell $t+1$, exactly the
first $t$ bits of $x$ copied onto the work tape (every other work cell still
blank), and work head at cell $t$.
\end{lemma}

\begin{lemma}[Transition from copy to rewind]
\lean{Complexity.pal_copy_end}
\label{Complexity.pal_copy_end}
\leanok
\uses{Complexity.palCfg, Complexity.palTM, Turing.Action.apply, Turing.Cfg.inputSymbol, Turing.MultiTapeTM.step, Turing.moveInputPos_neg_of_ne_left}
\difficulty{4}
Let $x$ be a binary word of length $n$. From the configuration of the three-state
palindrome machine on input $x$ in copy state with the input head on the right
boundary blank (cell $n+1$), all $n$ bits copied, and work head at cell $n$, one
transition switches to the rewind state and moves the input head left to cell
$n$, leaving the work tape and the work head unchanged.
\end{lemma}

\begin{lemma}[Single rewind-phase transition]
\lean{Complexity.pal_rewind_step}
\label{Complexity.pal_rewind_step}
\leanok
\uses{Complexity.palCfg, Complexity.palTM, Turing.Action.apply, Turing.Cfg.inputSymbol, Turing.MultiTapeTM.step, Turing.inputSymbolInner, Turing.moveInputPos, Turing.moveInputPos_neg_of_ne_left}
\difficulty{4}
Let $x$ be a binary word of length $n$ and $i < n$. From the configuration of the
three-state palindrome machine on input $x$ in rewind state with input head at
cell $i+1$, all $n$ bits copied, and work head at cell $n$, one transition moves
the input head left to cell $i$ and changes nothing else: the state, the copied
work tape, and the work head position are preserved.
\end{lemma}

\begin{lemma}[Rewind-phase invariant]
\lean{Complexity.pal_rewind}
\label{Complexity.pal_rewind}
\leanok
\uses{Complexity.palCfg, Complexity.palTM, Complexity.pal_rewind_step, Turing.MultiTapeTM.runFrom, Turing.MultiTapeTM.runFrom_succ_eq_step}
\difficulty{4}
Let $x$ be a binary word of length $n$. For every $i \le n$, running the
three-state palindrome machine for $i$ steps from the rewind-state configuration
on input $x$ with input head at cell $i$, all $n$ bits copied, and work head at
cell $n$ brings the input head to the left boundary cell $0$, with the state,
work tape, and work head unchanged.
\end{lemma}

\begin{lemma}[Transition from rewind to test]
\lean{Complexity.pal_test_start}
\label{Complexity.pal_test_start}
\leanok
\uses{Complexity.palCfg, Complexity.palTM, Turing.Action.apply, Turing.Cfg.inputSymbol, Turing.MultiTapeTM.step, Turing.moveInputPos, Turing.moveInputPos_pos_of_ne_right}
\difficulty{4}
Let $x$ be a binary word of length $n$. From the configuration of the three-state
palindrome machine on input $x$ in rewind state with the input head on the left
boundary blank (cell $0$), all $n$ bits copied, and work head at cell $n$, one
transition switches to the test state, moves the input head right to cell $1$,
and moves the work head left to the integer cell $n-1$ (which is $-1$ when $x$ is
empty), leaving the work tape unchanged.
\end{lemma}

\begin{lemma}[Work-tape read in the test phase]
\lean{Complexity.pal_test_read}
\label{Complexity.pal_test_read}
\leanok
\uses{Complexity.palCfg, Complexity.palTape, Turing.Cfg.workTapeSymbols}
\difficulty{3}
Let $x$ be a binary word of length $n$ and $i < n$. In the test-state
configuration of the palindrome machine on input $x$ with input head at cell
$i+1$, all $n$ bits copied, and work head at the integer cell $n-1-i$, the symbol
under the work head is the $i$-th bit of the reversal $x^{R}$.
\end{lemma}

\begin{lemma}[Single matching test-phase transition]
\lean{Complexity.pal_test_step}
\label{Complexity.pal_test_step}
\leanok
\uses{Complexity.palCfg, Complexity.palTM, Complexity.pal_test_read, Turing.Action.apply, Turing.Cfg.inputSymbol, Turing.MultiTapeTM.step, Turing.inputSymbolInner, Turing.moveInputPos, Turing.moveInputPos_pos_of_ne_right}
\difficulty{4}
Let $x$ be a binary word of length $n$ and $i < n$, and suppose the $i$-th bit of
the reversal $x^{R}$ equals the $i$-th bit of $x$. From the test-state
configuration of the palindrome machine on input $x$ with input head at cell
$i+1$, all $n$ bits copied, and work head at the integer cell $n-1-i$, one
transition advances the opposing heads: the input head moves right to cell $i+2$
and the work head moves left to cell $n-1-(i+1)$, with state and work tape
unchanged.
\end{lemma}

\begin{definition}[Agreement of opposing pairs from a position]
\lean{Complexity.palMatches}
\label{Complexity.palMatches}
\leanok
For a binary word $x$ and an index $i$, the predicate asserting that every
opposing pair at or after position $i$ agrees: for all $j$ with
$i \le j < |x|$, the $j$-th bit of the reversal $x^{R}$ equals the $j$-th bit
of $x$.
\end{definition}

\begin{lemma}[Peeling one matched pair off the agreement predicate]
\lean{Complexity.palMatches_succ}
\label{Complexity.palMatches_succ}
\leanok
\uses{Complexity.palMatches}
\difficulty{4}
Let $x$ be a binary word and $i < |x|$ an index at which the $i$-th bit of the
reversal $x^{R}$ equals the $i$-th bit of $x$. Then all opposing pairs of $x$ at
or after position $i$ agree if and only if all opposing pairs at or after
position $i+1$ agree.
\end{lemma}

\begin{lemma}[Test-phase invariant]
\lean{Complexity.pal_test}
\label{Complexity.pal_test}
\leanok
\uses{Complexity.palCfg, Complexity.palMatches, Complexity.palMatches_succ, Complexity.palTM, Complexity.pal_test_read, Complexity.pal_test_step, Turing.Action.apply, Turing.Cfg.inputSymbol, Turing.MultiTapeTM.runFrom, Turing.MultiTapeTM.runFrom_of_halt, Turing.MultiTapeTM.runFrom_succ_eq_step, Turing.MultiTapeTM.runFrom_zero, Turing.MultiTapeTM.step, Turing.inputSymbolInner}
\difficulty{5}
Let $x$ be a binary word of length $n$, and let $r$ and $i \le n$ satisfy
$n = i + r$. Running the three-state palindrome machine for $r+1$ steps from the
test-state configuration on input $x$ with input head at cell $i+1$, all $n$ bits
copied, and work head at the integer cell $n-1-i$ produces a halted
configuration whose output is the single bit \emph{true} if every opposing pair
of $x$ at or after position $i$ agrees (the $j$-th bit of $x^{R}$ equals the
$j$-th bit of $x$ for all $i \le j < n$), and \emph{false} otherwise.
\end{lemma}

\begin{theorem}[Palindromes are decidable in linear time]
\lean{Complexity.PAL_mem_DTIME_linear}
\label{Complexity.PAL_mem_DTIME_linear}
\leanok
\uses{Complexity.DTIME, Complexity.PAL, Complexity.palCfg, Complexity.palMatches, Complexity.palTM, Complexity.pal_copy, Complexity.pal_copy_end, Complexity.pal_rewind, Complexity.pal_test, Complexity.pal_test_start, Turing.MultiTapeTM.indicator, Turing.MultiTapeTM.initCfg, Turing.MultiTapeTM.runFrom, Turing.MultiTapeTM.runFrom_add, Turing.MultiTapeTM.runFrom_succ_eq_step'}
\difficulty{5}
The language of binary palindromes — the set of words $x \in \{0,1\}^*$ with
$x^{R} = x$ — lies in $\mathrm{DTIME}(n+1)$: there exist a constant $c$ and a
multi-tape Turing machine over the binary alphabet that, on every input of length
$n$, halts within $c\,(n+1)$ steps and outputs true exactly when the input is a
palindrome, and false otherwise.
\end{theorem}

\begin{theorem}[Palindromes are decidable in polynomial time]
\lean{Complexity.PAL_mem_P}
\label{Complexity.PAL_mem_P}
\leanok
\uses{Complexity.P, Complexity.PAL, Complexity.PAL_mem_DTIME_linear, Complexity.mem_P_of_dtime_le}
\difficulty{1}
The language of binary palindromes — the set of words $x \in \{0,1\}^*$ equal to
their own reversal — belongs to the class $\mathsf{P}$ of languages decidable by
a multi-tape Turing machine in time polynomial in the input length.
\end{theorem}
